% TOP_PROBLEM: {"id": 30006750, "title": "Residual Finiteness of Hyperbolic Groups", "category_id": 17, "category": {"id": 17, "name": "group_theory", "display_name": "Group Theory", "description": "Problems about groups, group actions, representations, and related algebraic structures.", "slug": "group-theory", "order_index": 17, "created_at": "2026-07-31T15:26:25.671Z"}, "status": "open"}
% FIRST_POSTED: null
% !TeX program = lualatex
% Compile twice: lualatex 30006750.tex
% All references and prose are embedded; no BibTeX database or external inputs.
\documentclass[11pt,a4paper]{article}
\usepackage[margin=24mm,headheight=14pt]{geometry}
\usepackage{fontspec}
\setmainfont{Latin Modern Roman}
\setsansfont{Latin Modern Sans}
\setmonofont{Latin Modern Mono}
\usepackage{amsmath,amssymb,mathtools}
\usepackage{unicode-math}
\setmathfont{Latin Modern Math}
\usepackage{microtype}
\usepackage{enumitem,xurl,needspace}
\usepackage[dvipsnames]{xcolor}
\usepackage{hyperref}
\hypersetup{unicode=true,colorlinks=true,linkcolor=MidnightBlue,urlcolor=MidnightBlue,
 pdftitle={Research attempt: Problem 30006750},pdfauthor={Research notebook},bookmarksnumbered=true}
\usepackage{bookmark}
\usepackage{tocloft}
\setlength{\cftsecnumwidth}{3em}
\cftsetpnumwidth{2.5em}
\cftsetrmarg{4em}
\usepackage{titlesec}
\titleformat{\section}{\Large\bfseries\raggedright}{\thesection}{0.7em}{}
\usepackage{fancyhdr}
\pagestyle{fancy}\fancyhf{}
\fancyhead[L]{\small\sffamily Open Problems Math | Research notebook}
\fancyhead[R]{\small Problem 30006750 | 27 September 2026}
\fancyfoot[C]{\thepage}
\setlength{\parindent}{0pt}
\setlength{\parskip}{3pt plus 1pt}
\setlength{\emergencystretch}{3em}
\setcounter{tocdepth}{1}
\setcounter{secnumdepth}{1}
\setlist[itemize]{leftmargin=3em,itemsep=5pt,topsep=4pt,parsep=0pt}
\allowdisplaybreaks
\newcommand{\N}{\mathbb{N}}
\newcommand{\Z}{\mathbb{Z}}
\newcommand{\Q}{\mathbb{Q}}
\newcommand{\R}{\mathbb{R}}
\newcommand{\C}{\mathbb{C}}
\newcommand{\F}{\mathbb{F}}
\newcommand{\A}{\mathbb{A}}
\newcommand{\PP}{\mathbb P}
\newcommand{\E}{\mathbb{E}}
\newcommand{\eps}{\varepsilon}
\newcommand{\dd}{\,\mathrm d}
\newcommand{\sref}[2]{\hyperlink{source:#1:#2}{[S#2]}}
\newcommand{\eref}[2]{\hyperlink{extra:#1:#2}{[E#2]}}
% Turn the next switch off for a shorter reading copy. The quotations preserve
% every exactTarget field, including qualifications omitted from a short summary.
\newif\ifshowcatalogue
\showcataloguetrue
\newcommand{\cataloguescope}[1]{\ifshowcatalogue
 \paragraph{Exact scope in the supplied catalogue.}
 \begin{quote}\small #1\end{quote}\fi}
\usepackage{etoolbox}
\AtBeginEnvironment{itemize}{\small}
\numberwithin{equation}{section}
% Standalone export. Original research content preserved.
\begin{document}
\setcounter{section}{0}
% ================================================================
% problemId: 30006750
% Canonical problem ID: 30006750
% exactTarget SHA-256: f1c0589af9aa2a07145370f45da84b9cbc0359661ffcb73cbfc5c8df798a010e
\clearpage
\section*{\texorpdfstring{Residual Finiteness of Hyperbolic Groups}{Residual Finiteness of Hyperbolic Groups}}\label{30006750}
\subsection{Definitions and mathematical statement}
A group $G$ is residually finite if its finite quotients separate its elements from the identity:
\[
 \forall g\in G\setminus\{1\}\ \exists F\text{ finite}\ \exists\varphi:G\to F:
 \varphi(g)\ne1.
\]
The conjecture asks for this property for every word-hyperbolic group. A different quotient may be used for each $g$, and no algorithm or quantitative bound on $|F|$ is required. Hyperbolicity is defined using a finite generating set and the thin-triangle condition on its Cayley graph.
\par\smallskip\textit{Formulation and concept references:} \sref{30006750}{1}.
\cataloguescope{For every word-hyperbolic group G and every nonidentity element g in G, prove that there is a finite group F and a homomorphism \ensuremath{\varphi}:G\ensuremath{\to}F with \ensuremath{\varphi}(g)\ensuremath{\ne}1.}
\subsection{Short English statement}
Can every nonidentity element of every hyperbolic group be detected in some finite quotient?
% BEGIN RESEARCH ADDITION 143
\subsection{Research attempt}

\paragraph{Current frontier.}
The supplied paper constructs residually finite hyperbolic groups that are
not linear, so nonlinearity is not an obstruction to residual finiteness.
\sref{30006750}{1} Alekseev--Thom's 2026 manuscript obtains local embeddability
in finite groups for Kazhdan groups with a specified sofic embedding and an
ergodic centralizer action. It does not furnish those hypotheses for every
hyperbolic group. \eref{30006750}{2} Exact local embeddings must be distinguished
from approximate permutation actions. For finitely presented groups the
former condition is already equivalent to residual finiteness, as explained
in Bradford, Section 2.3. \eref{30006750}{1}

\paragraph{Attack route.}
Complete partial generator actions on a large Cayley ball to permutations
of a finite set. This works for free groups without an arithmetic
representation. For a hyperbolic presentation, the missing step is exact
simultaneous repair of relators without collapsing a chosen word. We prove
the precise local criterion and test both completion and a counterexample
strategy based on a concrete non-residually-finite one-relator group.

\paragraph{Attempt: a finite-presentation reduction.}
Let $G=\langle S\mid\mathcal R\rangle$ with finite $\mathcal R$, relator
lengths at most $L$, and symmetric word balls $B_R$. Suppose a finite group
$F$ admits an injection $\iota:B_R\to F$ satisfying
\[
 \iota(xy)=\iota(x)\iota(y)\quad(x,y,xy\in B_R).
\]
For $R\ge\max\{L,|g|,1\}$ the generator assignment extends to a homomorphism
$G\to F$ detecting $g\ne1$. Indeed $\iota(1)=1$ by idempotence, and
inverse pairs are respected. All prefixes of every relator lie in $B_R$,
so repeated local multiplication evaluates the relator to one. The
presentation gives the homomorphism, and the same prefix calculation sends
$g$ to $\iota(g)\ne1$.

Conversely, residual finiteness supplies a finite quotient detecting every
nonidentity element of the finite set $B_R^{-1}B_R$: take a product of the
individual finite quotients. This quotient is injective on $B_R$. Thus exact
local embeddings for every $R$ are equivalent to the target for a fixed
hyperbolic group, which is finitely presented. This is a self-contained
verification of the standard reduction, not a weaker conjecture solved
in place of residual finiteness. \eref{30006750}{1}

\paragraph{Free completion and its relator defect.}
For $F_d$, take the finite set of reduced words $B_R$. Right multiplication
$v\mapsto vs$ is a partial bijection where both endpoints lie in $B_R$.
Complete it to a permutation $P_s$ for each positive generator; unmatched
domains and ranges have equal size. Set $P_{s^{-1}}=P_s^{-1}$. With the
usual right-action convention these define an action of the free group,
because there are no relations to check. Every reduced word $w$ of length
at most $R$, read from the identity, follows its original path and ends at
$w$. The resulting finite permutation quotient detects all such words and
has order at most $|B_R|!$. The companion program constructs and checks
these actions on every word through its tested radius.

Arbitrary completion fails in the presence of relations. Even for the
radius-one ball $\{0,\pm x,\pm y\}$ of $\Z^2$, the completions
\[
 P_x=(-x\ 0\ x)(y\ {-y}),\qquad
 P_y=(-y\ 0\ y)(x\ {-x})
\]
retain the prescribed partial translations but do not commute:
$P_xP_y(0)=-y$ and $P_yP_x(0)=-x$, using rightmost-first composition in
this displayed calculation. Enlarging the ball forces relators only on
paths remaining in the interior, not on the newly completed boundary.

Nor is the boundary necessarily negligible. For $d\ge2$,
\[
 |B_R|=1+2d\frac{(2d-1)^R-1}{2d-2},\qquad
 \frac{|B_R\setminus B_{R-L}|}{|B_R|}
 \longrightarrow1-(2d-1)^{-L}.
\]
A fixed-thickness boundary can have positive limiting proportion. This
invalidates a particular truncation estimate, not soficity or residual
finiteness of these groups.

The exact repair statement needed is: for each hyperbolic presentation and
nontrivial $g$, find finite permutations satisfying every defining relator
identically while moving some basepoint along $g$. Hyperbolicity controls
fillings of infinite Cayley-complex loops; no passage from that fact to the
required exact finite permutation identities is proved here. Small Hamming
error would not meet this requirement.

\paragraph{A failed counterexample transfer.}
Consider
\[
 B=\langle a,b\mid ba^2b^{-1}=a^3\rangle,
 \qquad w=[bab^{-1},a].
\]
Every map from $B$ to a finite group kills $w$. If the image of $a$ has
order $m$, equality of orders of $a^2$ and $a^3$ forces
$\gcd(m,2)=\gcd(m,3)=1$. Therefore both powers generate $\langle a\rangle$
and conjugation by $b$ normalizes this cyclic group. Thus $bab^{-1}$
commutes with $a$.

But $w=bab^{-1}ab a^{-1}b^{-1}a^{-1}$ is nontrivial by HNN normal form:
none of its potential pinches has exponent divisible by the required
$2$ or $3$. The base element $a$ has infinite order. This is the standard
Baumslag--Solitar normal-form framework recalled in \eref{30006750}{3}.
It does not give a hyperbolic counterexample. In a word-hyperbolic group an
infinite-order element has positive stable translation length, invariant
under conjugacy and satisfying $\tau(a^k)=|k|\tau(a)$. The relation would
force $2\tau(a)=3\tau(a)$, impossible. Thus the very relation producing
this finite-quotient obstruction violates the loxodromic-element property
of hyperbolic groups.

A hyperbolic group merely surjecting onto $B$ does not inherit this failure.
The free group on $a,b$ already surjects onto $B$, while its residual
finiteness was proved by the finite-ball construction. A quotient map or
Rips-type surjection cannot by itself transfer the obstruction to its source.

\paragraph{Stress test.}
Finite presentation enters at the relator-prefix step; no equivalence of
LEF and residual finiteness is asserted for arbitrary infinitely presented
groups. Exact identities are essential, not a vanishing fraction of
violations. The example $B$ supplies an explicit nontrivial element killed
by every finite quotient, but fails hyperbolicity. The 2026 centralizer
criterion is useful only after its special sofic embedding and ergodicity
are established; neither is assumed for the target groups. \eref{30006750}{2}

\paragraph{Outcome.}
\textbf{Outcome: Exploratory approach with unresolved gap.}
A finite-ball separation construction for free groups and the exact obstruction
to its naïve extension to presentations have been supplied. A concrete
Baumslag--Solitar counterexample transfer has also been excluded. No new
residual-finiteness theorem is claimed. The remaining step is relation-preserving
finite completion, or construction of a genuinely hyperbolic group with an
element invisible in all finite quotients.

% END RESEARCH ADDITION 143
\subsection{Sources}
\begin{itemize}\raggedright
\item[\textnormal{[S1]}]\hypertarget{source:30006750:1}{} Nicolas Tholozan and Konstantinos Tsouvalas, “Residually Finite Non-Linear Hyperbolic Groups,” Commentarii Mathematici Helvetici 100 (2025), 1–9, DOI 10.4171/CMH/581.
\url{https://ems.press/content/serial-article-files/49988}.
{\small\textit{Catalogue formulation source.}}
% BEGIN RESEARCH REFERENCES 143
\item[\textnormal{[E1]}]\hypertarget{extra:30006750:1}{} H. Bradford,
\emph{Quantifying local embeddings into finite groups}, Journal of Algebra
608 (2022), 214--238; arXiv:2104.07111, Section 2.3.
\url{https://arxiv.org/html/2104.07111v2}.
\item[\textnormal{[E2]}]\hypertarget{extra:30006750:2}{} V. Alekseev and A. Thom,
\emph{Centralizers of sofic approximations of Kazhdan groups},
arXiv:2608.05362, 5 August 2026; ergodic-centralizer criterion.
\url{https://arxiv.org/abs/2608.05362}.
\item[\textnormal{[E3]}]\hypertarget{extra:30006750:3}{} M. Clay, M. Forester
and J. Louwsma, \emph{Stable commutator length in Baumslag--Solitar groups
and quasimorphisms for tree actions}, arXiv:1310.3861, Section 2.
\url{https://arxiv.org/pdf/1310.3861}.

% END RESEARCH REFERENCES 143
\end{itemize}

\end{document}
